\section{Structured physical OOD and the contraction interconnection}
\label{sec:gen-theory}
This appendix specifies the model and operating conditions behind
Appendix~\ref{sec:composite-extension}, then derives the prediction-only comparison.
The general composite controller is a theoretical extension; the reported
VLA runs retain the sampled laws in Section~\ref{sec:vla-method}.

\subsection{Physical matching and the same-command reference}
The frozen VLA supplies a command signal $r(t)$, which can be
piecewise continuous with locally finitely many jumps and can depend
causally on deployment observations.
Fix that entire supplied signal when defining the execution comparison.
The state $\xi$ contains the robot and necessary low-level controller
state; $k_0(\xi,r,t)$ is its nominal execution feedback. For manipulation,
include evolving object and contact/environment state needed to close these
dynamics, or explicitly bound its omitted influence below. Static geometry
can be a fixed common parameter. Endogenous object motion is not a common
input merely because both executions receive the same commands
(Appendix~\ref{app:physical-closure}). Write
\begin{equation}
 \dot\xi=f_0(\xi,t)+G(\xi,t)\mathcal A(\xi,u;\theta^*(t))+d(t),
 \qquad F_{\rm exec}=f_0+Gk_0(\xi,r,t).
 \label{eq:gen-plant}
\end{equation}
Controller dynamics are included in $f_0,G$ as needed. An allocator
chooses $u\in\mathcal U$ using $\hat\theta$. Assume
$\mathcal A(\xi,u;\hat\theta)-\mathcal A(\xi,u;\theta^*)=
\Psi(\xi,u,t)(\hat\theta-\theta^*)$ on the operating domain.
With $\delta_A=\mathcal A(\xi,u;\hat\theta)-k_0$,
$B_\theta=G\Psi$ and $d_{\rm eff}=d+G\delta_A$ give exactly
\eqref{eq:gen-matching}. In a reduced-state model, $d(t)$ is a pathwise remainder that can depend
on omitted physical state. A disjoint decomposition is
\[
\begin{aligned}
 d_{\rm eff}&=d_{\rm base}+d_{\rm hist}+d_{\rm contact}+G\delta_A,\\
 \|d_{\rm eff}\|_{M_c}&\le
 \bar d_{\rm base,M}+\bar d_{\rm hist,M}+
 \bar d_{\rm contact,M}+\bar d_{\rm real,M}.
\end{aligned}
\]
These bounds must hold uniformly over admissible corrected states,
commands, and omitted states within each certified smooth mode. Effects
already represented by matching or realization error are not counted again
as contact disturbance. Contact remainder bounds are assumptions, not
consequences of servo stability; impacts require the separate hybrid jump
conditions in Appendix~\ref{app:history-partition}, not a finite continuous
force bound. For gain and bias,
$\mathcal A=\operatorname{diag}(g)u+b$, the regressor contains command
and identity columns. Estimated gains bounded away from zero and feasible
commands are necessary for an inverse. Payload and friction changes need
a valid parameterization or bounded remainder. A locked joint can remove
authority. Perception changes can change the supplied commands; execution
tracking alone does not establish that those commands remain task-correct.

The nominal reference solves
$\dot\xi^{\rm nom}=F_{\rm exec}(\xi^{\rm nom},r(t),t)$ with a declared initial state.
It is the nominal physical response to the \emph{same commands actually
supplied during deployment}, not a counterfactual rollout that recomputes
VLA commands from a healthy observation stream. The reference may be an
analytical comparison system for prediction-only adaptation. Computing it
online is required only for a controller that uses its tracking signal.
For a reduced state, this is the response of the declared nominal reduced
model; a full-contact nominal simulation is a different comparison unless
its projection realizes that model or its discrepancy is bounded.

For every allowed command and each certified smooth operating mode,
assume $F_{\rm exec}$ is continuously differentiable in $\xi$, and $M_c(\xi,t)$ is $C^2$ in $\xi$ and $C^1$
in $t$ locally. Uniform positive metric bounds and
\eqref{eq:contraction-metric} hold on a stated operating domain $\Omega$
uniformly over $r\in\mathcal R$. The derivative $A_{\rm exec}$ holds $r$
fixed: it includes the robot and low-level feedback sensitivity, with no
VLA derivative. Although deployment commands can arise from feedback,
the nominal comparison and all virtual trajectories in this proof use
the same supplied $r(t)$. This is an input-conditioned execution
certificate, not a contraction certificate for the complete VLA loop.
Admissibility of the actual command signal and domain containment are
separate assumptions; the certificate does not establish them. Shared
commands do not imply shared object states or contact modes. A hybrid
comparison must cover the realized mode sequences and event timing; a
smooth servo metric alone supplies only its stated servo-domain guarantee.

Take a neighborhood on which every relevant pair has a minimizing
constant-speed geodesic contained in $\Omega$, and the squared metric
distance is smooth. A strongly geodesically convex normal neighborhood
is a sufficient local setting. The flow perturbations in first variation
must remain in $\Omega$ for sufficiently small time increments. Bounds
hold only until these conditions cease. Nonunique minimizers require a
separate upper-Dini argument; we do not assume a differentiable geodesic
choice through a cut locus.

The metric is independent of the command, the estimated parameter, and
the unknown parameter. Dependence on these quantities introduces extra
derivative or switching terms and requires a different certificate.
Similarly, $(\xi-\xi^{\rm nom})^\top M_c(\xi,t)(\xi-\xi^{\rm nom})/2$ is generally not
the geodesic half-energy. With a command-independent metric, a jump of
$r$ does not itself jump the energy if the physical/reference states are
continuous; the flow bound holds on each command interval and extends
across the jump by continuity. Direct reference resets, command-dependent
tracking-coordinate resets, and state or metric resets need explicit
jump bounds or a feasible interpolated reference. Physical controller
memory, actuator delays, and contact resets must be covered by the state
partition and sampled/hybrid conditions in
Appendix~\ref{app:history-partition}. VLA internal policy memory is not
part of the primary execution certificate.

\subsection{Identification information and the available feedback signal}
For measured or filtered regressions, declare weights $R_j\succ0$,
$\omega_j\ge0$, and form
\begin{align}
 y_j&=\Phi_j\theta^*(t_j)+\varepsilon_j,\qquad
 H_I=\sum_j\omega_j\Phi_j^\top R_j^{-1}\Phi_j,\nonumber\\
 b_I&=\sum_j\omega_j\Phi_j^\top R_j^{-1}y_j,\qquad
 \epsilon_I=b_I-H_I\theta^*(t).
 \label{eq:gen-information}
\end{align}
Then $b_I-H_I\hat\theta=-H_I\tilde\theta+\epsilon_I$ is the negative
gradient of the weighted prediction loss. Current data or a retained
informative dataset can supply this statistic. A command-history FIR
buffer supplies predictor context; it is not, by itself, an informative
identification memory. Estimation requires observability of the physical
parameter directions and a bias reference separating nominal mismatch
from the change being identified.

The geodesic signal $\sigma_{\rm exec}=B_\theta^\top p_\xi$ additionally
requires the nominal physical reference, realization sensitivity, and
metric endpoint covector. Appendix~\ref{app:gen-innerloop} constructs
these quantities from robot/controller dynamics. A low-level reference
generated from the current VLA chunk can be used when it realizes the
same nominal execution dynamics, or its dynamics defect is retained.
Neither the metric nor this tracking signal requires differentiating the
VLA. The evaluated prediction-only method computes neither quantity.
For a computed $\hat\sigma_{\rm exec}=\sigma_{\rm exec}+\epsilon_\sigma$ with
$\|\epsilon_\sigma\|\le\bar\epsilon_\sigma$, the parameter forcing in
\eqref{eq:gen-forcing} becomes
$k_c\sqrt{\gamma_+}\bar\epsilon_I+
\kappa\sqrt{\gamma_+}\bar\epsilon_\sigma+\nu/\sqrt{\gamma_-}$.
Approximate metrics or geodesics must justify this error bound and any
metric-inequality defect. If an approximate reference is used only to
compute $\hat\sigma_{\rm exec}$, its error is relative to the exact $\xi^{\rm nom}$ in the
theorem. Tracking a different reference satisfying
$\dot\xi^{\rm nom}=F_{\rm exec}(\xi^{\rm nom},r,t)+\delta_r$ adds the separate first-variation term
$-p_r^\top\delta_r$, where $p_r=M_c(\xi^{\rm nom},t)c_s(0)$. Its magnitude is
at most $X\|\delta_r\|_{M_c(\xi^{\rm nom},t)}$; this reference-dynamics defect
must be added to the tracking forcing $\bar d_M$. Approximate quantities
cannot simply be substituted as exact ones.

\subsection{Prediction-only contraction and feedback coupling}
\label{sec:prediction-comparison}
With $\kappa=0$, \eqref{eq:gen-law} is prediction-gradient adaptation in
the parameter metric $\Gamma^{-1}$. Under a common exogenous stream
$(H_I,b_I)$, two unprojected observers have variation
\begin{equation}
 \delta\dot{\hat\theta}=-k_c\Gamma H_I\delta\hat\theta,\qquad
 \frac{d}{dt}\frac12\|\delta\hat\theta\|_{\Gamma^{-1}}^2
 =-k_c\delta\hat\theta^\top H_I\delta\hat\theta.
 \label{eq:observer-differential}
\end{equation}
If $H_I\succeq\mu I$, this is conditional observer contraction at rate
$k_c\mu\gamma_-$. The projected case follows from normal-cone
monotonicity (Appendix~\ref{app:gen-proof}), without differentiating
through active faces. On the robot, correction changes the state and
therefore its data stream; this conditional observer statement does not
certify the feedback interconnection.

To account for that dependence, suppose throughout the operating domain
\begin{align}
 \|M_c^{1/2}B_\theta\Gamma^{1/2}\|&\le L_B,\qquad
 \|d_{\rm eff}\|_{M_c}\le d_0+\ell_xX+\ell_zZ,\nonumber\\
 \|\epsilon_I\|&\le\epsilon_0+m_xX+m_zZ,
 \qquad \|\dot\theta^*\|\le\nu.
 \label{eq:gen-coupling}
\end{align}
All constants are nonnegative. For this execution/estimation comparison
\eqref{eq:prediction-comparison}, define
\begin{align}
 a&=\lambda-\ell_x,& b&=L_B+\ell_z,\nonumber\\
 c&=k_c\mu\gamma_- -k_c\sqrt{\gamma_+}m_z,&
 d&=k_c\sqrt{\gamma_+}m_x,\qquad
 q_z=k_c\sqrt{\gamma_+}\epsilon_0+\nu/\sqrt{\gamma_-}.
 \label{eq:prediction-coefficients}
\end{align}
\begin{proposition}[Prediction-only interconnection]
\label{prop:prediction-comparison}
Assume the geometric, matching, projection, and information conditions
above on an interval beginning at $t_e$, with $H_I\succeq\mu I$.
The $\kappa=0$ law obeys \eqref{eq:prediction-comparison}. Write
$\mathcal A_c=\left[\begin{smallmatrix}-a&b\\d&-c\end{smallmatrix}\right]$
and $r_0=(d_0,q_z)^\top$. Componentwise,
\begin{equation}
 \begin{bmatrix}X(t)\\Z(t)\end{bmatrix}\le
 e^{\mathcal A_c(t-t_e)}\begin{bmatrix}X(t_e)\\Z(t_e)\end{bmatrix}
 +\int_{t_e}^t e^{\mathcal A_c(t-s)}r_0\,ds.
 \label{eq:prediction-envelope}
\end{equation}
If $a,c>0$ and $ac>bd$, this comparison matrix is Hurwitz, and
\begin{equation}
 \limsup_{t\to\infty}X(t)\le\frac{cd_0+bq_z}{ac-bd},\qquad
 \limsup_{t\to\infty}Z(t)\le\frac{dd_0+aq_z}{ac-bd},
 \label{eq:prediction-floor}
\end{equation}
provided all assumptions remain valid for that infinite interval.
\end{proposition}
These are asymptotic comparison bounds, not startup bounds. A sufficient
invariant rectangle contains $(X(t_e),Z(t_e))$ and satisfies
$a\bar X\ge b\bar Z+d_0$, $c\bar Z\ge d\bar X+q_z$; its whole
trajectory tube must remain in the domain. Alternatively propagate
\eqref{eq:prediction-envelope} directly. Information acquisition requires
its own transient bound. If $d=0$, both diagonal rates being positive
suffice even for large $b$, although the transient may be large.
The sampled coupled-error result in Appendix~\ref{app:smallgain} expresses
the same feedback structure with the actual estimator and application schedule.

\subsection{Recovery, regime changes, and task scope}
For Theorem~\ref{thm:gen-adaptive}, if
$W(t_e)>W_{\rm tar}>Q/\beta$, a sufficient additional recovery time is
\begin{equation}
 T_{\rm rec}=\beta^{-1}\log
 \frac{W(t_e)-Q/\beta}{W_{\rm tar}-Q/\beta}.
 \label{eq:gen-time}
\end{equation}
Information-acquisition time must be added and its tube enforced. A bound
$\max\{W(t_e),Q/\beta\}<R$ proves containment only if the entire
radius-$R$ tube about the same-command response lies in the certified
domain and command admissibility persists. This is \emph{execution
recovery}: approach to the nominal physical response under the supplied
commands, up to estimation, model, timing, and realization errors.
If a separate task specification guarantees success throughout
$\|\xi-\xi^{\rm nom}\|<r_{\rm task}$ around a task-valid nominal response, one may
take $R\le\sqrt{\underline m}r_{\rm task}$. Both path and terminal
conditions must be included. A healthy VLA rerun generally supplies a
different command sequence; comparing it requires an additional argument
for command divergence or for the complete policy/plant interconnection.

For example, $\dot z=-z+r$ contracts at rate one for a common supplied
$r(t)$ in the metric $M_c=1$. Closing the feedback $r=2z$ instead gives
$\dot z=z$. Execution contraction therefore proves neither full-loop
stability nor task success. Dropping a VLA derivative while claiming
contraction of that full loop would be invalid. No such derivative or
full-loop contraction assumption enters the primary execution theorem.

A parameter jump is distinct from bounded-rate drift. If the physical and
reference states, metric, $\Gamma$, and estimate are continuous or unchanged,
\begin{equation}
 W(t_j^+)\le W(t_j^-)+\|\Delta\theta^*(t_j)\|_{\Gamma^{-1}}.
 \label{eq:gen-jump}
\end{equation}
Restart the envelope with valid post-change information. A large $H_I$
may encode obsolete parameters; Appendix~\ref{app:gen-regimes} bounds
stale-data bias, reacquisition, and repeated changes. Separate reference,
state, or metric resets require their own jump bounds. None of these
tracking-storage results proves pairwise contraction of every augmented
$(\xi,\hat\theta)$ trajectory. That stronger claim needs a differential
metric inequality for the actual augmented dynamics, including
state-dependent regressors, normalizers, and metric derivatives.
